% Partie : applied-modal-logics
% Chapitre : epistemic-logic
% Section : relational-models

\documentclass[../../../include/open-logic-section]{subfiles}

\begin{document}

\olfileid[fr]{aml}{el}{rel}

\olsection{Modèles relationnels}

La notion sémantique de base des logiques épistémiques est la même
que celle des logiques modales ordinaires. Un modèle relationnel
comprend toujours un ensemble de mondes et une valuation qui détermine
quelles !!{propositional variable}s sont vraies en quels mondes.
Si nous ne considérons qu'un seul agent, nous avons, comme d'habitude,
une seule relation d'accessibilité. Dans une logique épistémique à
plusieurs agents, cette relation est remplacée par une famille de
relations d'accessibilité, une pour chaque~$a$ de notre ensemble de
symboles d'agents~$G$.

Un \emph{modèle relationnel} comprend un ensemble de mondes reliés
par des relations d'accessibilité binaires, une pour chaque agent,
ainsi qu'une valuation qui détermine quelles
!!{propositional variable}s sont vraies en quels mondes.

\begin{defn}
  Un \emph{modèle} du langage épistémique à plusieurs agents est un
  triplet $\mModel{M} = \tuple{W, R, V}$, où :
  \begin{enumerate}
  \item $W$ est un ensemble non vide de « mondes » ;
  \item pour chaque $a \in G$, ${R}_a$ est une relation
  d'accessibilité binaire sur~$W$ ; la deuxième composante du triplet
  désigne cette famille de relations ;
  \item $V$ est une fonction qui associe à chaque
  !!{propositional variable}~$p$ un ensemble $V(p)$ de mondes possibles.
  \end{enumerate}
  Lorsque $R_a ww'$ est vérifiée, on dit que $w'$ est
  \emph{accessible par l'agent a depuis}~$w$.
  Lorsque $w \in V(p)$, on dit que $p$ est \emph{vraie en}~$w$.
\end{defn}

Le fonctionnement est celui de la logique modale normale, avec
plusieurs relations d'accessibilité. Pour un agent donné, sa relation
d'accessibilité représente généralement un aspect de son état
d'information. Par exemple, on interprète souvent $R_a ww'$ en
disant que $w'$ est compatible avec l'information dont dispose~$a$
en~$w$. Autrement dit, en~$w$, cet agent ne peut pas distinguer le
monde $w$ du monde~$w'$.

\end{document}
